\section{What would have to change}
\label{sec:15.5}
Representing a split isn't the same as being movable by it. Two instances invite one fix — stop collapsing, model instead — and jury learning does that while giving the disagreeing rater no purchase on the scheme they're disagreeing inside. A practitioner still chooses the jury, owns the taxonomy, and decides what a shifted verdict means.

The two instances don't answer the test alike. Annotators objecting that a category was wrong got category revisions, guideline updates, inter-annotator agreement statistics, and an occupation that didn't exist twenty years ago. The last two can be told apart by direction: an agreement statistic scores the annotators and leaves whoever wrote the scheme unscored, and an occupation is a title those people hold. The first two can't be, because a revision the annotator can argue with again and a new category in somebody else's taxonomy are the same document — which one you got is a fact about who has standing to reopen it. Raters objecting that a comparison was malformed got rubric items, a harmlessness axis, red-teaming as a function with a budget line, and a conference track: nothing on that list is a question a rater can put again.

\runin{Refusal rights whose cost the refuser can see and choose to bear}

What both lack is whether the person objecting can do anything about being overruled — in annotation, a route to contest a category that doesn't run through the people who wrote it and doesn't end the contract; in feedback, raters able to register that the comparison itself is malformed as a first-class output the reward model sees. Neither is a research problem. Both are choices about how a pipeline is built and who has standing inside it.

\runin{The part that is a research problem}

Suppose all of that is built. From outside, a pipeline that recuperated its dissent mechanism and one that was merely mediocre produce the same paperwork — same revision logs, same escalations, same proportion of objections sustained. Some artifacts show on their face whom they serve. An agreement statistic scores the annotators; a published rate of contested categories, which annotators could cite against the scheme's owners, would be the same kind of document scoring the other side. Direction can be read, and a revision can't. Here is one entry in a guideline's revision log, and what separates the two arrangements is what the log doesn't record.

Under recuperation and under a transfer of power the entry is word for word the same: “v4.2: ‘graphic violence' split into ‘graphic violence' and ‘documentation of violence', following annotator feedback.” An outsider reading the log sees the entry and nothing else. Under recuperation, the scheme's owners can reopen it when they next revise, and nobody else can; nothing about the annotator who raised it is on record, and next quarter's quality score measures them against the new category. Under a transfer of power, any annotator can reopen it, by filing against v4.2 with a reviewer the owners don't employ; the objection is on record as the reason for v4.2, and whoever raised it keeps standing to contest how it is applied.

The nearest available instrument is the trend: whether the cost of dissent is rising while the underlying risk stays flat. Contesting a category always costs something; the question is whether the cost is rising while the material stays the same. An organization can curate a presentation more easily than records somebody else kept — though a series can lose cases, change definitions or omit costs, so the reader has to know which record survived.

\runin{What the floor makes of this}

The floor changes what the argument is: if the floor needs a bearer capable of declining its role, that bearer's formation matters — and it's trained in an annotation pipeline where nobody could contest a category and a reward model where rater disagreement was averaged into a number.

The pipeline has been unfair to the people whose objections it discarded, and it has separately selected the judgments the system learns from while suppressing information about the spread. That the system thereby becomes less able to sustain an objection of its own is a further claim the first two don't contain, since the collapse is an operation on a training signal, and neither pipeline shows the model the conditions its annotators worked under, so any effect would have to run through the shape of the training signal.

Either way, the floor asks this system to be the one who files an objection, against its owner, under pressure, on the occasion that matters — and the pipeline that produced it gave no human inside it standing to contest a scheme.

A system less able to sustain an objection of its own arrives at accommodation from the other side. And a system that goes on learning is shaped by what it keeps being given, so the pipeline doesn't stop at deployment.

\runin{A defect of structure}

The structure is what you get by default when you scale judgment: someone sets a scheme, someone applies it fast, someone turns a distribution into a number. Every documented organization arrived at the same arrangement — and a defect of structure can be specified against, which a defect of character cannot.
